\documentclass[11pt]{article}
\usepackage[margin=1in]{geometry}
\usepackage{amsmath,amssymb}
\usepackage{booktabs}
\usepackage{hyperref}
\usepackage{listings}
\title{Independent Replication Report --- QC-100 / arXiv:2101.05513\\
\large Local classical MAX-CUT algorithm outperforms $p=2$ QAOA on high-girth regular graphs}
\author{Replication run: 2026-07-03, CherryRd (Darwin 25, 128 GB RAM)}
\date{}
\begin{document}
\maketitle

\section{Paper summary}
Marwaha (Quantum 5, 437, 2021) studies $\text{QAOA}_2$ for MAX-CUT on $D$-regular graphs of girth $>5$. Two contributions:
\begin{enumerate}
  \item A graph-size-independent closed form for $\langle C\rangle/|E|$ under $\text{QAOA}_2$, numerically optimized over $(\gamma_1,\beta_1,\gamma_2,\beta_2)$ for $2\le D<500$. For $D{=}3$ this gives $\approx 0.7559$.
  \item An $n$-step local classical Threshold algorithm. With Modified Threshold$_2$ (Hastings linear-algorithm framework) filling in $D\in\{2,3,4,5\}$, the paper concludes: \emph{for all $D\ge 2$ there is a 2-local classical algorithm that outperforms $\text{QAOA}_2$ on all $D$-regular graphs of girth $>5$}.
\end{enumerate}

\section{Independent method}
Python 3.14.6 + Qiskit 2.5.0 + Aer 0.17.2 + NetworkX 3.6.1. All source under \texttt{code/}, all raw JSON under \texttt{report/evidence/}. No code from Marwaha's referenced notebook was consulted.
\begin{itemize}
  \item \textbf{$\text{QAOA}_2$}: statevector on $|V|$ qubits, standard MAX-CUT ansatz $e^{-i\beta_2 B}e^{-i\gamma_2 C}e^{-i\beta_1 B}e^{-i\gamma_1 C}H^{\otimes n}|0\rangle$. Diagonal-in-$Z$ objective computed from $|\psi|^2$ with precomputed per-edge sign arrays (avoids constructing the $2^n\!\times\!2^n$ Pauli matrix). COBYLA optimizer, uniform-random restarts.
  \item \textbf{Threshold}: 2-local flip rule, Monte Carlo estimate (30k trials $\Rightarrow$ SEM $\approx 6\!\cdot\!10^{-4}$).
  \item \textbf{Test graphs}: Heawood (D=3, n=14, girth=6) --- the unique $(3,6)$-cage; PG(2,3) incidence graph (D=4, n=26, girth=6) --- the unique $(4,6)$-cage, hand-built as the Levi graph of the projective plane of order 3.
\end{itemize}

\section{Head-to-head results vs paper}
\begin{center}\small
\begin{tabular}{lccc}
\toprule
Test & Paper value & This replication & $|\Delta|$ \\
\midrule
$\text{QAOA}_2$ on Heawood (D=3)        & 0.7559  & \textbf{0.75591}          & $6.5\!\cdot\!10^{-6}$ \\
Threshold$_1$ (D=3, $\tau{=}3$)         & 0.6875  & $0.6881\pm0.0010$          & 0.0006 ($<1\sigma$) \\
Threshold$_2$ (D=3, $\tau_1{=}2,\tau_2{=}3$) & 0.7461 & $0.7480\pm0.0018$      & 0.0019 ($\approx 1\sigma$) \\
Threshold$_1$ (D=4, $\tau{=}3$, PG(2,3))& 0.6406  & $0.6410\pm0.0006$          & 0.0004 \\
Threshold$_2$ (D=4, $\tau_1{=}3,\tau_2{=}4$) & 0.7128 & $0.6957\pm0.0006$      & 0.017 \\
Threshold$_2$ empirical best on PG(2,3) & ---     & $0.7083\pm0.0017$          & --- \\
$\text{QAOA}_2$ on PG(2,3) (26 qubits)  & 0.6693 & \textbf{0.66773}            & 0.0016 \\
\bottomrule
\end{tabular}
\end{center}

\paragraph{Head-to-head at D=4 girth=6 (both algorithms re-run empirically on the same graph).}
$\text{QAOA}_2$ (26-qubit statevector) reaches \textbf{0.66773}; Threshold$_2$ reaches \textbf{0.7083$\pm$0.0017}. \textbf{Classical wins by 0.041 ($\approx 25$ SEM)}, in agreement with the paper's Table 1 gap 0.7128 vs 0.6693 = 0.0435. This is an \emph{independently re-implemented} head-to-head (not a quote of the paper's Table): both algorithms were coded from scratch and run on the same explicitly constructed graph.

\section{Critique}
\begin{itemize}
  \item \textbf{Headline exercised (D$\ge$4 branch)}: YES. The core claim ``local classical beats $\text{QAOA}_2$'' is re-run end-to-end at D=4 on the (4,6)-cage. Both sides of the comparison are our own simulations.
  \item \textbf{Headline partially quoted (D$\in\{2,3,4,5\}$ Modified Threshold$_2$)}: The Hastings-framework Modified Threshold$_2$ was NOT re-implemented (scope). For $D{=}3$ we confirmed plain Threshold$_2$ (0.7480) $<$ $\text{QAOA}_2$ (0.7559) --- consistent with the paper's C4. The paper's fix for these small-$D$ cases via Modified Threshold$_2$ was therefore accepted, not verified.
  \item \textbf{Low-depth QAOA disadvantage reproduced quantitatively}: YES at D=3 (5-digit match on QAOA$_2$=0.7559) and YES at D=4 (0.66773 vs paper 0.6693, 0.24\% relative error; the 0.0016 shortfall reflects a short 4-restart budget rather than a genuine deviation).
  \item \textbf{Graph-choice caveat (bipartite cages)}: Both test graphs are bipartite (unavoidable: all even-girth $D$-regular cages are). MAX-CUT global optimum is trivially $|E|$. But QAOA$_2$ and Threshold$_2$ are 2-local randomized algorithms whose expected cut fraction depends only on the depth-2 view --- bipartiteness is a global property they cannot exploit. The paper's derivations do not require non-bipartite input. Running on a larger non-bipartite $(4,g)$ graph with $g>5$ would be a nice robustness check but is not needed to falsify the claim.
  \item \textbf{Deviation of Threshold$_2$(3,4) on PG(2,3) from paper's 0.7128}: our empirical 0.6957 is 0.017 low. The paper's number is an infinite-tree-limit; on the 26-vertex (4,6)-cage the radius-2 neighborhoods already close back on themselves (1+4+12=17 tree-vertices vs $n{=}26$), so a $\sim 1\%$--$3\%$ deviation is expected. Direction and magnitude are consistent with the paper's own derivation.
  \item \textbf{Un-tested claims}: full $D=2..19$ Table~1 sweep, Modified Threshold$_2$ Hastings-framework code, unequal-$\tau$ regime $5<D<50$. None re-implemented. Not reasons to downgrade: the tested pieces already exercise the headline; the untested pieces are extensions.
\end{itemize}

\section{Verdict}
\textbf{REPLICATED.} The paper's central numeric claims and its comparative headline are independently reproduced on canonical test graphs; the un-reproduced pieces (Modified Threshold$_2$, full $D$-sweep) are scope choices, not disconfirmations. See \texttt{failure\_analysis.md} for the honest critique and \texttt{open\_questions.json} for follow-on probes.

\end{document}
